\section{本章纲要：失地后的国家怎样组织江海}

南宋不是北宋的缩小版，也不只是金与蒙古之间等待覆亡的政权。行在、两浙、江汉、四川、两淮与港口必须共同供给军队、官僚和迁徙人群；金朝的边界政策、海上贸易和地方市场又不断改变这些区域能够承担的成本。南渡以后的核心问题，是如何把东南的粮、盐、货币、船只和文书转成防务，同时不把不同地方的生活经验压缩成“富庶江南”。

本章从绍兴、隆兴等和议与战事入手，但不把外交当作唯一主线。会子、漕运、军屯、市舶、契约、寺院和宗族，是国家与地方相互进入的不同接口；襄阳、临安和崖山则显示战争在区域间以不同节奏推进。宋代政书与官修史须同契约、碑刻、港口遗存和地方材料并读，才能分开财政规范、实际执行和家庭承受。
